\chapter{Theoretical Background}
\label{ch:theoretical_background}

\section{Thermal Control in Space Engineering}
In aerospace engineering, the Thermal Control System (TCS) is a mission-critical subsystem responsible 
for maintaining all spacecraft components within their allowable temperature ranges during all mission phases \cite{tcs_slides0, tfg_carol}. 
Unlike terrestrial applications, spacecraft operate in a harsh and highly variable thermal environment. 
Electronic components, batteries, structural joints, and sensitive optical payloads must be protected 
against severe thermal gradients and extreme temperature fluctuations \cite{spacethctrl, tfg_carol}. Failure 
to manage these temperatures can lead to accelerated material degradation, mechanical distortion, or 
catastrophic failure of the payload. % \cite{tfg_carol}. 

The primary objectives of the \acrshort{tcs} are: to ensure optimal performance during operational modes
and to guarantee survival during non-operational or safe modes \cite{tcs_slides0}. This is achieved through
a careful balance of internal heat dissipation and external environmental fluxes, utilizing 
both passive and active thermal control technologies \cite{tcs_slides0,th_handbook}.

\section{The Space Thermal Environment}
The thermal environment in space is radically different from that on Earth. The vacuum 
of space precludes any convective heat transfer to the ambient environment, leaving thermal radiation 
as the sole mechanism for rejecting heat to deep space, which acts as a blackbody sink at 
approximately $3\text{ K}$ \cite{tfg_carol,tcs_slides0}.

Spacecraft, particularly those in Low Earth Orbit (LEO), are subjected to
 three primary sources of external environmental heating \cite{tfg_carol}:
\begin{itemize}
    \item \textbf{Direct Solar Radiation:} The most significant heat source, characterized by the solar constant (roughly $1361 \text{ W/m}^2$ at $1\text{ AU}$).
    \item \textbf{Earth Albedo:} The fraction of solar radiation reflected by the Earth's surface and atmosphere. It is highly dependent on orbit altitude, attitude, and local weather conditions.
    \item \textbf{Planetary Infrared Radiation:} The thermal energy emitted by the Earth itself, typically ranging between $200$ and $300 \text{ W/m}^2$ for LEO altitudes.
\end{itemize}

During a typical orbit, a spacecraft transitions between periods of direct sunlight (eclipse out) and 
Earth's shadow (eclipse in), generating severe thermal cycling \cite{tfg_carol}. In contrast, Deep Space 
or Geostationary Earth Orbit (GEO) missions face different challenges, such as prolonged unilateral solar 
exposure or deep cold soaking %\cite{tfg_carol}.

\section{Heat Transfer Mechanisms}
Although convection is generally negligible in the vacuum of space (except within pressurized crewed 
modules or actively pumped fluid loops) \cite{tcs_slides4}, conduction and radiation are the foundational 
mechanisms governing spacecraft thermal dynamics %\cite{a_pdf}.

\subsection{Conduction Heat Transfer}
Conduction dominates the internal heat distribution through the spacecraft's solid structures. 
It is governed by Fourier's Law of Heat Conduction, which in its one-dimensional macroscopic form is 
expressed as \cite{tcs_slides1, hmt_f&a}:
\begin{equation}
    \dot{Q} = -k A \frac{dT}{dx}
\end{equation}
where $\dot{Q}$ is the heat transfer rate ($\text{W}$), $k$ is the thermal conductivity of the material
($\text{W/m}\cdot\text{K}$), $A$ is the cross-sectional area ($\text{m}^2$), and $dT/dx$ is the temperature
gradient.

\subsubsection{Thermal Contact Resistance and Thermal Straps}
An added complexity in space structures is the Thermal Contact Resistance (TCR). When two solid 
surfaces are bolted or pressed together, microscopic roughness limits the actual contact area, creating 
resistance to heat flow \cite{tcs_slides3}. In the vacuum of space, the absence of interstitial air drastically
increases this resistance. To mitigate this or to intentionally decouple components, thermal engineers use 
interstitial fillers (e.g., Cho-Therm, Sigraflex). Alternatively, highly conductive flexible 
links known as \textit{Thermal Straps} (typically made of copper or pyrolytic graphite) are used to transport 
heat from sensitive electronics to radiators while mechanically decoupling them to prevent vibration 
transmission \cite{tcs_slides3, tcs_slides0}.

\subsection{Radiation Heat Transfer}
Thermal radiation is the emission of electromagnetic waves from all matter at temperatures above absolute zero. 
It is the only way a spacecraft can reject waste heat \cite{heattransfer}. 

\subsubsection{Blackbody Radiation}
An ideal blackbody is a theoretical entity that absorbs all incident radiation and emits the maximum 
possible thermal energy. The total emissive power is given by the Stefan-Boltzmann Law \cite{heattransfer}:
\begin{equation}
    E_b = \sigma T^4
\end{equation}
where $\sigma \approx 5.67 \times 10^{-8} \text{ W/m}^2\cdot\text{K}^4$ is the Stefan-Boltzmann constant. 
The spectral distribution of this energy is described by Planck's Law, and the peak emission wavelength shifts 
inversely with temperature according to Wien's Displacement Law ($\lambda_{\text{max}} T \approx 2897 \times 10^{-3} \text{ m}\cdot\text{K}$) %\cite{a_pdf}.

\subsubsection{Thermo-Optical Properties of Real Surfaces}
Real surfaces deviate from blackbody behavior. Their thermal interaction is governed by:
\begin{itemize}
    \item \textbf{Absorptivity ($\alpha$):} The fraction of incident radiation absorbed. In spacecraft design, solar absorptivity ($\alpha_s$) is a critical parameter.
    \item \textbf{Emissivity ($\epsilon$):} The ratio of energy radiated by a material to that radiated by a blackbody at the same temperature. Infrared emissivity ($\epsilon_{IR}$) dictates the heat rejection capability.
\end{itemize}
The $\alpha_s / \epsilon_{IR}$ ratio is the fundamental driver of a surface's equilibrium temperature in space. For example, Optical Solar Reflectors (OSRs) feature a very low $\alpha_s$ and high $\epsilon_{IR}$, making them ideal for radiators \cite{tcs_obj}.

\subsubsection{Radiative Exchange and View Factors}
The net radiative heat transfer between two diffuse, gray surfaces depends heavily on their 
geometric orientation, quantified by the View Factor (or configuration factor) $F_{ij}$ \cite{heattransfer}.
This factor represents the fraction of radiation leaving surface $i$ that directly intercepts surface $j$. 
The reciprocity theorem states that $A_i F_{ij} = A_j F_{ji}$, which is extensively used to simplify complex 
thermal network calculations.

\begin{figure}[htpb]
    \centering
    % \includegraphics[width=0.6\textwidth]{figures/view_factor_diagram.png}
    \caption{Geometric representation of the view factor $F_{ij}$ between two differential surfaces.}
    \label{fig:view_factor}
\end{figure}

\section{Mathematical Modeling of Thermal Systems}
To accurately predict spacecraft temperatures, Thermal Mathematical Models (TMM) must be constructed 
and solved %\cite{tcs_obj, a_pdf}.

\subsection{Limitations of the Finite Element Method (FEM)}
While the Finite Element Method (FEM) is the industry standard for structural analysis, 
it is often computationally prohibitive for system-level space thermal analysis. FEM relies heavily on 
dense continuous meshes which are excellent for conduction but struggle with the highly non-linear, 
non-local nature of radiative heat exchange and the calculation of thousands of view factors %\cite{a_pdf}.

\subsection{Lumped Parameter Method (LPM)}
The aerospace industry overwhelmingly relies on the Lumped Parameter Method (LPM), which discretizes
 the spacecraft into a network of isothermal thermal nodes %\cite{a_pdf, tcs_obj}. 
 The transient energy balance for a generic node $i$ is formulated as an ordinary differential equation:
\begin{equation}
    C_i \frac{dT_i}{dt} = \sum_{j \neq i} L_{ij} (T_j - T_i) + \sum_{j \neq i} R_{ij} (T_j^4 - T_i^4) + Q_{i}
\end{equation}
where:
\begin{itemize}[label={\scriptsize\raisebox{0.5ex}{\textbullet}}]
  \item$C_i$ [J/K]: thermal capacitance of node $i$ ($m_i\ c_p,_i$)
  \item$T_i$ [K]: temperature of node $i$
  \item$t$ [s]: time
  \item$L_{ij}$ [W/K]: conductive coupling (linear conductance)
  \item$R_{ij}$ [$W/K^4$]: radiative coupling factor (incorporating view factors and emissivities)
  \item$Q_i$ [W]: net heat source applied to node $i$ (internal dissipation and absorbed external fluxes)
\end{itemize} 

This approach provides a perfect balance between computational efficiency and the precision required for 
complex radiative environments %\cite{a_pdf}.

\section{Advanced Thermal Control Technologies}
Translating the mathematical model into reality requires specific hardware. A robust TCS relies on \cite{tcs_slides0,th_handbook}:
\begin{itemize}
    \item \textbf{Multi-Layer Insulation (MLI):} Blankets made of multiple layers of reflective films (e.g., Kapton, Mylar) separated by low-conductance netting. They drastically reduce radiative heat transfer, effectively isolating the spacecraft from the environment.
    \item \textbf{Heat Pipes:} Two-phase heat transfer devices that use the latent heat of vaporization of a working fluid (usually ammonia in space) to transport large amounts of heat across the spacecraft with a near-zero temperature gradient.
    \item \textbf{Survival Heaters:} Resistive electrical patch heaters, often governed by mechanical thermostats or thermistors, ensuring that sensitive components (like batteries and propellant tanks) do not drop below their survival limits during eclipse phases.
\end{itemize}

\section{Software and Advanced Simulation Methods}
System-level thermal analysis is performed using specialized computational tools. In the European space sector,
\textbf{ESATAN-TMS} is the standard software suite. It solves the nodal networks defined by the LPM and 
calculates radiative view factors and both steady-state and transient thermal models %\cite{a_pdf}. 
Similarly, \textbf{Thermal Desktop} is widely used in NASA ecosystems %\cite{a_pdf}.

\subsection{Post-Processing and Python Integration}
While ESATAN-TMS is a powerful solver, its native visualization capabilities can be rigid. %\cite{a_pdf}. 
Modern thermal engineering workflows increasingly integrate programming languages like \textbf{Python} to 
automate model execution, extract raw nodal data, and generate advanced analytics %\cite{a_pdf}. 
Using libraries such as \texttt{NumPy} and \texttt{Panda}, engineers can process massive datasets, 
visualize transient temperature curves across hundreds of nodes, plot heat flow networks, and execute rapid 
parametric optimization to refine the thermal design %\cite{a_pdf}.


